\documentclass[11pt]{report}
\usepackage[a4paper,margin=2.7cm]{geometry}
\usepackage{amsmath,amssymb,amsthm}
\usepackage{booktabs}
\usepackage{physics}
\usepackage{graphicx}
\usepackage{hyperref}
\usepackage[round]{natbib}
\usepackage{siunitx}
\usepackage{xcolor}
\usepackage{tcolorbox}
\tcbuselibrary{skins,breakable}

\newtcolorbox{checkbox}[1][]{colback=blue!4!white,colframe=blue!40!black,
  boxrule=0.4pt,arc=1mm,left=6pt,right=6pt,top=4pt,bottom=4pt,breakable,
  title={Check},fonttitle=\bfseries\small,coltitle=blue!30!black,#1}
\newtcolorbox{choicebox}[1][]{colback=orange!5!white,colframe=orange!45!black,
  boxrule=0.4pt,arc=1mm,left=6pt,right=6pt,top=4pt,bottom=4pt,breakable,
  title={Choice made},fonttitle=\bfseries\small,coltitle=orange!30!black,#1}

\newcommand{\code}[1]{\texttt{#1}}
\newcommand{\Msun}{M_{\odot}}

\title{\Large Wave-Optics Gravitational Lensing of Gravitational Waves\\
\large in a Hierarchical Triple: a Worked Derivation}
\author{Final project --- Gravitational Waves and AI-Assisted Research}
\date{September 2026}

\begin{document}
\maketitle

\begin{abstract}
\noindent
This chapter derives, from the linearized wave equation, the formalism used
to compute how a gravitational wave (GW) from a compact binary is lensed by
a third, more massive body in the same hierarchical triple. It follows the
scalar-field reduction of GW propagation in a weak-field lens background
through to the Kirchhoff--Fresnel diffraction integral, works out the
point-mass lens in closed form and cross-checks it against two independent
numerical methods, derives the specific geometry of a source and lens that
share one physical system (as opposed to the usual well-separated
cosmological lens), and closes with the two regimes --- a static lens during
a fast chirp, and a moving lens sampled by a slowly-evolving quasi
-monochromatic source --- that \code{cases/case\_A\_chirp} and
\code{cases/case\_B\_monochromatic} instantiate numerically. Every numbered
check below is implemented and run in \code{tests/}; this document derives
\emph{why} each check is the right one, not just that it passes.
\end{abstract}

\tableofcontents

% ======================================================================
\chapter{Setup and scope}
% ======================================================================

\section{The physical picture}

The system is a hierarchical triple: an inner compact binary of masses
$m_1,m_2$ (the GW \emph{source}) orbits, at semi-major axis $a_\mathrm{in}$,
a common centre of mass which itself orbits a third, more massive body of
mass $M_L$ (the \emph{lens}) at semi-major axis $a_\mathrm{out}$, with
\begin{equation}
  a_\mathrm{in} \ll a_\mathrm{out}.
  \label{eq:hierarchy}
\end{equation}
This is the ordinary stability condition for a hierarchical triple, and it
does double duty here: it is also exactly the condition under which the
inner binary can be treated, for lensing purposes, as a single point source
--- both members are deflected identically, so there is no differential
lensing across the binary itself (no "binary microlensing" substructure).
Throughout, $M_L \gg m_1+m_2$, so the lens sits, to the accuracy used here,
at rest at the triple's barycentre while the inner binary (as a whole)
orbits it.

The triple is taken to be within our own Galaxy (or a similarly nearby
system), so the lens and source are at essentially the same distance from
Earth and cosmological redshift plays no role: distances $D_L$ (observer to
lens), $D_S$ (observer to source) and $D_{LS}$ (lens to source) are ordinary
Euclidean distances, and the lens redshift is $z_L=0$ (the factor
$(1+z_L)$ that multiplies the lens mass everywhere in the cosmological
lensing literature is simply $1$ here). Chapter~\ref{ch:geometry} works out
a consequence of this that a textbook cosmological treatment does not have
to face: $D_{LS}$ is \emph{not} a large, slowly-varying distance --- it is
set by the outer orbit itself, and can be small, zero, or even formally
negative.

\subsection*{Two epochs of the same inspiral}
The inner binary loses energy to gravitational radiation and inspirals. Two
snapshots of that inspiral are studied:
\begin{description}
  \item[Case A --- the chirp.] Late inspiral: the GW frequency sweeps
    through the sensitive band of a ground-based-like detector in seconds.
    That is enormously faster than any plausible outer orbital period, so
    the lens is \emph{frozen} at a single impact parameter for the whole
    observation.
  \item[Case B --- quasi-monochromatic.] Early inspiral: the inner binary is
    wide, the GW frequency is low and drifts negligibly over an observation
    spanning several outer periods. Now the outer orbital motion \emph{is}
    resolved, and the impact parameter sweeps periodically.
\end{description}
Both regimes are made precise, and checked against the actual numbers used
in \code{cases/}, in Chapter~\ref{ch:regimes}.

\begin{choicebox}
The lens mass, $M_L\sim5\times10^4\,\Msun$ (an ``IMBH'', between the
conventional intermediate-mass and supermassive black hole ranges), is
chosen so that the wave-optics parameter $w$ (defined in
Eq.~\ref{eq:wdef}) is not astronomically extreme at either epoch's
frequency, not because it is a typical hierarchical-triple lens mass in the
literature. \code{wiki/log.md} records this choice and the numbers it
produces; it is not a fitted or tuned value beyond that one requirement.
\end{choicebox}

% ======================================================================
\chapter{From the wave equation to a scalar field}
\label{ch:scalar}
% ======================================================================

The user-facing goal of this project stops here: reaching the point where
gravitational-wave lensing is a \emph{scalar}-field diffraction problem, and
justifying why that is legitimate.

\section{The wave equation on a weak-field background}

Write the full metric as a background plus a small perturbation,
$g_{\mu\nu}=\bar g_{\mu\nu}+h_{\mu\nu}$, where $\bar g_{\mu\nu}$ is the
lens's own weak-field metric (Newtonian potential $U=-GM_L/r$, $|U|/c^2\ll1$
everywhere relevant here since the lens is far outside its own Schwarzschild
radius at the impact parameters considered) and $h_{\mu\nu}$ is the GW. In
Lorenz gauge on this background, the linearized vacuum Einstein equations
give a wave equation for $h_{\mu\nu}$ with a curvature-coupling term,
\begin{equation}
  \bar\Box h_{\mu\nu} - 2\bar R_{\mu\alpha\nu\beta}h^{\alpha\beta} = 0 .
  \label{eq:linearized}
\end{equation}
This is the standard starting point for GW propagation on a curved
background (e.g.\ \citealp{Isaacson1968}). The high-frequency (eikonal)
regime relevant here is set by comparing the GW wavelength $\lambda_{GW}$ to
the background's curvature length scale $\mathcal{R}$: the ray/thin-lens
treatment below is valid when $\lambda_{GW}\ll\mathcal{R}$. For a point lens
$\mathcal{R}\sim R_E$ (the Einstein radius), and the curvature-coupling term
in Eq.~\eqref{eq:linearized} is then subdominant to $\bar\Box h_{\mu\nu}$ by
a factor $\sim(\lambda_{GW}/\mathcal{R})^2$ (\citealp{Peters1974};
\citealp{Takahashi2003}, \S2). Dropping it, each independent polarization
amplitude $h_{\mu\nu}=e_{\mu\nu}\,\Phi$ (fixed polarization tensor
$e_{\mu\nu}$, transverse-traceless) obeys the \emph{same} equation as a
massless scalar field on the lens background,
\begin{equation}
  \bar\Box\,\Phi = 0 .
  \label{eq:scalar_wave}
\end{equation}
This is the key reduction: lensing in this regime does not mix or rotate
the two GW polarizations, and every polarization is amplified/retarded by
the identical (complex, frequency-dependent) factor derived below. It is
also why solving the scalar-field diffraction problem \emph{is} solving the
GW lensing problem here, and nothing about spin-2 structure survives past
this point.

\begin{checkbox}
This reduction is only used, never re-derived numerically in this project
--- it is a standard result cited above. What \emph{is} checked
numerically, repeatedly below, is everything built on top of it: the
diffraction integral it leads to, the point-lens closed form, and its
geometric-optics limit.
\end{checkbox}

\section{The Helmholtz equation and the thin lens}

For the weak-field metric $\bar g_{\mu\nu}=\eta_{\mu\nu}+2U\,\delta_{\mu\nu}$
(isotropic weak-field form, $c=1$ units in this section), Eq.~\eqref{eq:scalar_wave}
for a wave of fixed angular frequency $\omega$, $\Phi(\vb r,t)=\psi(\vb
r)e^{-i\omega t}$, becomes
\begin{equation}
  \left(\nabla^2 + \omega^2\right)\psi = 4\omega^2 U\,\psi ,
  \label{eq:helmholtz}
\end{equation}
a Helmholtz equation with the lens potential acting as a (frequency
-dependent) refractive-index perturbation, $n(\vb r)=1-2U(\vb r)$ ---
exactly analogous to geometric/wave optics in a medium of varying refractive
index, which is the origin of calling this whole formalism ``wave optics.''

\subsection*{Thin-lens (Fresnel--Kirchhoff) reduction}
Because $a_\mathrm{out}\ll D_L,D_S$ (Eq.~\ref{eq:hierarchy} and the Galactic
-scale distances above), the deflection happens in a region whose extent
along the line of sight is negligible compared to $D_L$ and $D_S$: the
standard thin-lens approximation. Propagating Eq.~\eqref{eq:helmholtz} from
the source plane to the lens plane and on to the observer via the Fresnel
diffraction integral (paraxial approximation, valid since all angles
involved are $\ll1$) gives the amplification factor --- the ratio of the
lensed to unlensed wave at the observer,
\begin{equation}
  F(f) \;=\; \frac{D_S}{D_L D_{LS}}\,\frac{f}{i}
  \int d^2\vb\xi\; \exp\!\left[\, i 2\pi f\, t_d(\vb\xi;\vb\eta)\,\right],
  \label{eq:Fdim}
\end{equation}
where $\vb\xi$ is position on the lens plane, $\vb\eta$ the (fixed) source
position, and $t_d$ the time delay (geometric $+$ Shapiro) of a ray through
$\vb\xi$ relative to the undeflected path
(\citealp{SchneiderEhlersFalco1992}; \citealp{Takahashi2003}). By
construction, $h_{\rm lensed}(f)=F(f)\,h_{\rm unlensed}(f)$: everything the
lens does to the (scalar-equivalent) waveform is contained in this single
complex, frequency-dependent number.

\subsection*{Dimensionless form}
Scaling the lens-plane position by the Einstein radius $\xi_0\equiv R_E$ and
the source position by the corresponding source-plane scale $\eta_0\equiv
\xi_0 D_S/D_L$, with $\vb x\equiv\vb\xi/\xi_0$, $\vb y\equiv\vb\eta/\eta_0$
(so $y=|\vb y|$ is the dimensionless impact parameter used throughout this
project), and defining the dimensionless deflection potential $\psi(\vb x)$
via $t_d = (1+z_L)\dfrac{D_S}{D_L D_{LS}}\left[\tfrac12|\vb x-\vb
y|^2-\psi(\vb x)\right]$, Eq.~\eqref{eq:Fdim} becomes
\begin{equation}
  F(w,\vb y) \;=\; \frac{w}{2\pi i}\int d^2\vb x\;
  \exp\!\Big[\,i\,w\,\varphi(\vb x,\vb y)\,\Big],
  \qquad
  \varphi(\vb x,\vb y)\equiv \tfrac12|\vb x-\vb y|^2-\psi(\vb x),
  \label{eq:Fdimless}
\end{equation}
where the dimensionless frequency is
\begin{equation}
  w \;\equiv\; 8\pi\,\frac{GM_L}{c^3}\,f
  \label{eq:wdef}
\end{equation}
(the $(1+z_L)$ factor is $1$ here, per Chapter~1). This is exactly the
function \code{F\_bruteforce\_2d} in \code{waveoptics.py}: no closed form
has been used yet, only Eq.~\eqref{eq:helmholtz} and the thin-lens/paraxial
reduction.

\begin{checkbox}[title={Check: $w\to0$}]
Physically, a wave much longer than the lens's effect scale should not be
lensed at all: $F(w,\vb y)\to1$ as $w\to0$. This is not obvious from
Eq.~\eqref{eq:Fdimless} by naive Taylor expansion (the integral is only
conditionally convergent, \S\ref{sec:pointlens}), but it is confirmed
numerically to $\sim5\times10^{-4}$ at $w=10^{-4}$ by
\code{tests/test\_waveoptics.py::check\_low\_w\_limit}.
\end{checkbox}

% ======================================================================
\chapter{The point-mass lens}
\label{ch:pointlens}
% ======================================================================
\label{sec:pointlens}

Everything above is general (any $\psi(\vb x)$). From here on the lens is a
single point mass, $\psi(\vb x)=\ln|\vb x|$, the deflection potential whose
gradient $\nabla\psi=\vb x/|\vb x|^2$ gives the standard $1/|\vb x|$
deflection law of a point mass.

\section{Geometric optics: images, magnification, time delay}
\label{sec:geo}

By symmetry take $\vb y=(y,0)$. Stationary points of $\varphi(\vb x,\vb y)$
(Fermat's principle: images sit where the time delay is stationary) restricted
to the $x_2=0$ axis solve
\begin{equation}
  x - \frac{1}{x} = y
  \;\;\Longrightarrow\;\;
  x_\pm = \tfrac12\left(y\pm\sqrt{y^2+4}\right),
  \qquad x_+x_-=-1 .
  \label{eq:images}
\end{equation}
Two images, always: $x_+>1>0>x_->-1$ for $y>0$.

\subsection*{Magnification}
For an axisymmetric lens the 2D lens map $\vb y(\vb x)=f(x)\hat{\vb x}$,
$f(x)\equiv x-1/x$, has Jacobian $\det(\partial\vb y/\partial\vb x) =
\big(f(x)/x\big)\,f'(x)$ (the product of the tangential and radial stretch
factors --- a standard identity for axisymmetric maps,
e.g.\ \citealp{SchneiderEhlersFalco1992}). With $f'(x)=1+1/x^2$,
\begin{equation}
  \mu(x) \;=\; \frac{1}{\det(\partial\vb y/\partial\vb x)}
  \;=\; \frac{1}{\left(1-\tfrac{1}{x^2}\right)\left(1+\tfrac{1}{x^2}\right)}
  \;=\; \frac{x^4}{x^4-1} .
  \label{eq:magnification}
\end{equation}

\begin{checkbox}[title={Check: total magnification vs.\ Paczy\'nski (1986)}]
The total (unsigned) magnification $\mu_{\rm tot}(y)=|\mu(x_+)|+|\mu(x_-)|$
built from Eqs.~\eqref{eq:images}--\eqref{eq:magnification} must equal the
independent, textbook microlensing formula
$A(y)=(y^2+2)/\!\left(y\sqrt{y^2+4}\right)$ \citep{Paczynski1986}. Verified
to floating-point precision ($<3\times10^{-15}$, all $y$ tested) in
\code{tests/test\_waveoptics.py::check\_paczynski\_magnification} --- a
genuine cross-check, since Eq.~\eqref{eq:magnification} was derived here
from the Jacobian, not copied from Paczy\'nski's paper.
\end{checkbox}

\subsection*{Time delay and the geometric-optics amplification factor}
The (dimensionless) time delay of each image is $\varphi(x_\pm,y)$, and the
delay of the (always positive-parity, ``minimum'') image $x_+$ relative
to itself is trivially zero, while the saddle-point image $x_-$ arrives
later by $\Delta T(y)\equiv\varphi(x_-,y)-\varphi(x_+,y)>0$. Evaluating
Eq.~\eqref{eq:Fdimless} by stationary phase (valid as $w\to\infty$) sums the
two images with Morse-index phases (minimum: $0$; saddle: $-\pi/2$, i.e.\ an
overall factor of $-i$, in the $h(f)=\int h(t)e^{-i2\pi ft}dt$ convention
fixed in \code{wiki/conventions.md}) \emph{and} each image's own,
un-subtracted phase $w\,\varphi(x_+,y)$:
\begin{equation}
  F_{\rm geo}(w,y) \;=\;
  e^{\,iw\varphi(x_+,y)}\left[\sqrt{\mu(x_+)} \;-\; i\sqrt{|\mu(x_-)|}\,
  e^{\,iw\Delta T(y)}\right].
  \label{eq:Fgeo}
\end{equation}

\begin{checkbox}[title={A mistake, and how it was caught}]
An earlier version of Eq.~\eqref{eq:Fgeo} omitted the leading
$e^{iw\varphi(x_+,y)}$ factor (implicitly assuming a reference phase
subtraction that Eq.~\eqref{eq:Fdimless} does not make). This gave the
right $|F|$ but the wrong phase compared to the exact closed form
(\S\ref{sec:closedform}) at large $w$. The error was diagnosed by
comparing the two directly and recognizing the residual phase scaled
linearly with $w$ at the rate $\varphi(x_+,y)$ --- and fixed by restoring
the missing factor. After the fix, \code{tests/test\_waveoptics.py::
check\_geometric\_optics\_limit} finds relative agreement
$3\times10^{-4}$ at $w=10$, shrinking to $3\times10^{-5}$ at $w=1000$, the
expected behaviour of an asymptotic ($w\to\infty$) expansion. See
\code{wiki/log.md} for the full account.
\end{checkbox}

\section{The exact result: a confluent hypergeometric function}
\label{sec:closedform}

Solving Eq.~\eqref{eq:Fdimless} for the point-mass lens exactly (originally
\citealp{Peters1974}; \citealp{Deguchi1986}; the standard form quoted in the
GW-lensing literature is \citealp{Takahashi2003}) gives
\begin{equation}
  F(w,y) \;=\;
  \exp\!\left[\frac{\pi w}{4} + i\frac{w}{2}\ln\frac{w}{2}\right]
  \Gamma\!\left(1-\frac{iw}{2}\right)
  {}_1F_1\!\left(\frac{iw}{2},\,1;\,\frac{iwy^2}{2}\right).
  \label{eq:Fclosed}
\end{equation}
This is \code{src/gwlens/waveoptics.py::F\_point\_lens}, evaluated with
arbitrary-precision complex arithmetic (\code{mpmath}, since
\code{scipy.special.hyp1f1} does not accept complex parameters). Rather than
trust Eq.~\eqref{eq:Fclosed} on citation alone, it is checked against two
independently-derived numerical evaluations of Eq.~\eqref{eq:Fdimless}
itself.

\subsection*{Independent check 1: the 1D Bessel-reduced integral}
Writing $|\vb x-\vb y|^2=x^2+y^2-2xy\cos\theta$ in polar coordinates on the
lens plane and using the standard identity
$\int_0^{2\pi}d\theta\,e^{ia\cos\theta}=2\pi J_0(a)$
(Abramowitz \& Stegun 9.1.21) to do the angular integral in
Eq.~\eqref{eq:Fdimless} analytically leaves a 1D radial integral,
\begin{equation}
  F(w,y) \;=\; -i\,w\,e^{\,iwy^2/2}\int_0^\infty dx\;x\,J_0(wxy)\,
  \exp\!\left[i w\left(\tfrac12 x^2-\ln x\right)\right].
  \label{eq:Fradial}
\end{equation}
This integral is only conditionally convergent; it is regularized with the
standard $i\epsilon$ prescription (\citealp{UlmerGoodman1995}),
$w\to w(1+i\eta)$, i.e.\ a convergence factor $e^{-\eta w x^2/2}$, evaluated
for a sequence of $\eta\to0$ and extrapolated. \code{src/gwlens/waveoptics.py::
F\_radial\_1d} implements this, with the cutoff radius and the number of
quadrature sub-intervals set automatically from $\eta$ and $w$ (so as to
resolve the local oscillation period rather than a hand-tuned grid).

\begin{checkbox}[title={Check: $O(\eta)$ convergence to the closed form}]
At $w=1,\,y=1$: relative error vs.\ Eq.~\eqref{eq:Fclosed} is
$2.6\times10^{-2}$ at $\eta=0.02$, $1.3\times10^{-2}$ at $\eta=0.01$,
$6.5\times10^{-3}$ at $\eta=0.005$ --- each halving of $\eta$ almost exactly
halves the error (ratios $1.98$, $1.99$), the expected $O(\eta)$ scaling of
the regularization itself. \code{tests/test\_waveoptics.py::
check\_radial\_1d\_converges\_to\_closed\_form}.
\end{checkbox}

\subsection*{Independent check 2: brute-force 2D quadrature}
Eq.~\eqref{eq:Fdimless} evaluated directly by 2D quadrature over a
regularized (Gaussian-damped) square grid, with no angular reduction and no
closed form anywhere --- the least clever, most first-principles evaluator
in the project. At $w=1,\,y=1$, $x_{\max}=25$, $1200\times1200$ grid points,
$\eta=0.005$: relative error $1.1\times10^{-2}$, consistent with the same
$O(\eta)$ regularization error seen in the 1D check.
\code{tests/test\_waveoptics.py::check\_bruteforce\_2d\_agrees};
\code{waveoptics.py::F\_bruteforce\_2d}.

\begin{checkbox}[title={Check: Einstein-ring diffraction spike}]
At exact alignment ($y=0$) geometric optics diverges ($\mu(x_+{=}1)$ is
singular, Eq.~\ref{eq:magnification}) --- the textbook point-source
caustic. The exact $F(w,0)$ from Eq.~\eqref{eq:Fclosed} stays finite for
every $w$ tested ($1.81$, $3.96$, $7.93$ at $w=1,5,20$): diffraction
regularizes the caustic, the GW analogue of the optical Arago/Poisson spot
(\citealp{Deguchi1986}). \code{tests/test\_waveoptics.py::
check\_einstein\_ring\_regularization}.
\end{checkbox}

\section{A fast hybrid evaluator}

Eq.~\eqref{eq:Fclosed}'s confluent hypergeometric series converges slowly
once its argument $iwy^2/2$ is large --- exactly the regime, large $w$,
where Eq.~\eqref{eq:Fgeo} is independently validated to agree with it to
better than $0.4\%$. \code{waveoptics.py::F\_hybrid} therefore uses
Eq.~\eqref{eq:Fclosed} below $w=30$ and Eq.~\eqref{eq:Fgeo} above it (the
threshold checked directly, \code{tests/test\_waveoptics.py::
check\_hybrid\_matches\_at\_threshold}: $3.5\%$ agreement at $w=30$ across
several $y$). This was not an optimization made in advance: evaluating
Eq.~\eqref{eq:Fclosed} directly across the frequency sweep actually used in
Case~A took $\sim150$\,s; \code{F\_hybrid} needs none of that (see
\code{wiki/log.md}).

% ======================================================================
\chapter{The hierarchical triple as a lens system}
\label{ch:geometry}
% ======================================================================

\section{The outer orbit}
The source's position relative to the (effectively fixed) lens follows an
ordinary Kepler orbit: elements $(a_\mathrm{out},e_\mathrm{out},
i_\mathrm{out},\Omega_\mathrm{out},\omega_\mathrm{out})$, solved by Newton's
method on Kepler's equation $M=E-e\sin E$ and rotated into the observer
frame ($z$ along the line of sight) by the standard sequence
$R_z(\Omega)R_x(i)R_z(\omega)$ (e.g.\ \citealp{MurrayDermott1999}):
\begin{align}
  x_{\rm sky} &= r\left[\cos\Omega\cos\theta - \sin\Omega\sin\theta\cos i\right],\\
  y_{\rm sky} &= r\left[\sin\Omega\cos\theta + \cos\Omega\sin\theta\cos i\right],\\
  z_{\rm los} &= r\sin\theta\sin i, \qquad \theta\equiv\nu+\omega,
  \label{eq:rotation}
\end{align}
$r,\nu$ the orbital radius and true anomaly.

\begin{checkbox}
Kepler's equation solver inverts to machine precision for
$e\in\{0,0.3,0.7,0.95\}$; a circular orbit gives exactly constant $r$; an
edge-on orbit ($i=\pi/2$, $\Omega=\omega=0$) projects to a line
($y_{\rm sky}\equiv0$ to $6\times10^{-17}$); a face-on orbit ($i=0$)
reproduces the orbital-plane position exactly. All in
\code{tests/test\_geometry.py}.
\end{checkbox}

\section{$D_{LS}$ is not what it is in a textbook lens}

This is the one place this project's geometry genuinely differs from the
cosmological lensing this formalism was built for, and it is easy to get
wrong by carrying over the textbook intuition unexamined. In the usual
setting, lens and source are at very different, essentially fixed distances
and $D_{LS}$ barely varies. Here they are in the \emph{same} system: the
lens sits at $D_L$, and the source's distance is
$D_S(t)=D_L+z_{\rm los}(t)$, with $z_{\rm los}(t)$ from
Eq.~\eqref{eq:rotation} at the AU-to-sub-pc scale of the outer orbit. Since
$|z_{\rm los}|/D_L\sim10^{-8}$ for the parameters used here, $D_S\approx D_L$
is an excellent approximation \emph{everywhere it appears except} in
\begin{equation}
  D_{LS}(t) = D_S(t) - D_L = z_{\rm los}(t)
\end{equation}
itself, where it is the whole story. The angular Einstein radius,
\begin{equation}
  \theta_E(t)^2 = \frac{4GM_L}{c^2}\,\frac{D_{LS}(t)}{D_L D_S(t)}
  \;\approx\; \frac{4GM_L}{c^2}\,\frac{z_{\rm los}(t)}{D_L^2},
  \label{eq:thetaE_t}
\end{equation}
is therefore itself time-dependent, and \textbf{undefined whenever
$z_{\rm los}(t)\le0$} --- the source is instantaneously nearer Earth than the
lens, so the lens is not between the observer and the source at all, and
there is no lensing geometry for that half of the orbit (not an
approximation artefact: $F=1$ there by construction). As
$z_{\rm los}(t)\to0^+$, $\theta_E\to0$ and $y=\theta/\theta_E(t)\to\infty$
for any nonzero transverse offset $\theta$: the formalism turns lensing off
smoothly and on its own at the front/back crossing, with no special-casing
needed beyond excluding $z_{\rm los}\le0$.

This is exactly the origin of \emph{repeated} lensing in a hierarchical
triple: only the far-side half of each outer orbit is lensed at all, so a
periodically-orbiting source produces periodic lensing pulses, once per
outer period, rather than continuous modulation
(\citealp{DOrazioLoeb2020}, in the geometric-optics limit; the same
structure appears here in wave optics, Case~B, \S\ref{sec:caseB}).

\begin{checkbox}
For the (near edge-on, $i_\mathrm{out}=87^\circ$) orbit used in
\code{cases/}, almost exactly half the orbit is lensed
($49.7$--$50.0\%$ across independent samplings), and $y$ is finite and
positive throughout the lensed half.
\code{tests/test\_geometry.py::check\_time\_varying\_impact\_parameter}.
\end{checkbox}

\section{Where the thin-lens/paraxial approximation itself is being pushed}
\label{sec:paraxial}

\S\ref{ch:scalar} derived Eq.~\eqref{eq:Fdim} under a paraxial (small-angle)
Fresnel approximation, which additionally requires the physical Einstein
radius on the lens plane, $\xi_0$, to be much smaller than the propagation
distance $D_{LS}$ itself --- an assumption that is automatic in cosmological
lensing (where $D_{LS}$ is a fixed, large distance) but is exactly what
Eq.~\eqref{eq:thetaE_t} shows is \emph{not} automatic here, since
$D_{LS}(t)$ shrinks to zero twice per outer orbit. Using
$\xi_0=\sqrt{4GM_L D_{LS}/c^2}$ (the point-lens Einstein radius written with
$D_S\approx D_L$, \S\ref{ch:geometry}), the paraxial condition
$D_{LS}\gg\xi_0$ reduces to
\begin{equation}
  D_{LS}(t) \;\gg\; \frac{4GM_L}{c^2} \;=\; 2\,R_{\rm Sch}(M_L) ,
  \label{eq:paraxial_condition}
\end{equation}
twice the lens's own Schwarzschild radius --- the \emph{same} weak-field
condition already invoked in \S\ref{ch:scalar} to drop the curvature
-coupling term in Eq.~\eqref{eq:linearized}, now applied to the
line-of-sight separation instead of the impact parameter. This is a
genuine, if narrow, respect in which this project's geometry pushes past
where the standard formalism was built to operate, and it is worth being
precise about exactly how much it matters, rather than either ignoring it
or re-deriving non-paraxial diffraction (a substantially harder problem,
evaluated below as not worth it here).

\begin{table}[h]
\centering
\begin{tabular}{lr}
\toprule
quantity & value (this triple) \\
\midrule
$R_{\rm Sch}(M_L)$ & $9.87\times10^{-4}$ AU \\
$D_{LS}$ at the deepest-lensing point ($y=y_{\min}$) & $1.81$ AU \\
$D_{LS}/(2R_{\rm Sch})$ there & $919$ \\
$\xi_0/D_{LS}$ there (direct paraxial-angle check) & $0.033$ \\
time per crossing spent with $D_{LS}\lesssim2R_{\rm Sch}$ & $60$ s \\
that time as a fraction of $P_{\rm out}$ & $1.7\times10^{-4}$ \\
\bottomrule
\end{tabular}
\caption{The paraxial condition, Eq.~\eqref{eq:paraxial_condition}, evaluated for the system in \code{src/gwlens/system.py} by \code{theory/paraxial\_validity.py} (\code{theory/paraxial\_validity\_numbers.json}).}
\label{tab:paraxial}
\end{table}

Table~\ref{tab:paraxial} is the answer: at the orbital phase that actually
produces the interesting physics in Case~B --- the point of closest
projected approach, $y=y_{\min}$, where the lensing pulses in
\code{cases/case\_B\_monochromatic/caseB\_repeated\_pulses.png} peak ---
$D_{LS}$ is close to its own \emph{maximum}
($\sim a_\mathrm{out}\sin i_\mathrm{out}$, since minimum transverse offset
and maximum line-of-sight separation occur near the same orbital phase for
this geometry), not its minimum, and exceeds $2R_{\rm Sch}$ by three orders
of magnitude. The paraxial approximation only becomes marginal in a window
of order $60$ seconds around each front/back crossing -- a fraction
$\sim1.7\times10^{-4}$ of the $4$-day outer period -- and that window is
\emph{already} the region where $y\to\infty$ and $F\to1$ by construction
(\S\ref{ch:geometry}), not a region this project draws any conclusion from.
A genuine non-paraxial (wide-angle) diffraction treatment near the crossing
would be needed to say anything quantitative about that $60$-second window
itself, and would require abandoning the Fresnel integral, Eq.~\eqref{eq:Fdim},
for a substantially more involved formulation -- a well-posed but
research-level extension, not attempted here. Given that it would only
refine an already-negligible, already-correctly-approximated ($F\approx1$)
sliver of the signal, the added complexity is judged not to be worth it for
this project's scope: this is a stated limitation, checked and quantified
rather than an unexamined one.

% ======================================================================
\chapter{The inner binary as a source}
% ======================================================================

The unlensed waveform uses the leading-order (Newtonian/quadrupole,
restricted) inspiral: not full post-Newtonian, since the physics under test
here is the lensing, not waveform accuracy (a stated scope choice, not an
oversight). The frequency evolves as
\begin{equation}
  \dot f = \frac{96}{5}\pi^{8/3}\left(\frac{G\mathcal{M}}{c^3}\right)^{5/3} f^{11/3},
  \qquad
  \mathcal{M}\equiv\frac{(m_1m_2)^{3/5}}{(m_1+m_2)^{1/5}},
  \label{eq:dfdt}
\end{equation}
the standard quadrupole-formula chirp rate (\citealp{Peters1964};
\citealp{PetersMathews1963}; \citealp{Maggiore2008}). Integrating
Eq.~\eqref{eq:dfdt} directly gives the time to (formal) coalescence and the
frequency track,
\begin{equation}
  \tau(f) = \frac{5}{256}(\pi f)^{-8/3}\left(\frac{G\mathcal{M}}{c^3}\right)^{-5/3},
  \qquad
  f(t) = \frac{1}{\pi}\left[\frac{5}{256}\left(\frac{G\mathcal M}{c^3}\right)^{-5/3}\right]^{-3/8}\!\!(t_c-t)^{-3/8}.
  \label{eq:ft}
\end{equation}

\begin{checkbox}[title={Check, and a bug caught by it}]
$f(t{=}0)$ from Eq.~\eqref{eq:ft} must reproduce the $f_0$ used to define
$t_c=\tau(f_0)$ (self-consistency of the inversion) and must agree with
direct numerical integration of Eq.~\eqref{eq:dfdt} (\code{scipy.integrate.
solve\_ivp}). An early hand-inversion of $\tau(f)$ for $f(t)$ mis-placed an
exponent (a bracket-power slip) and failed the first check immediately;
after fixing it, agreement with the ODE integration is $6\times10^{-9}$.
\code{tests/test\_chirp.py}.
\end{checkbox}

The (restricted, leading-order) time-domain strain amplitude,
\begin{equation}
  h(t) = \frac{4}{D_{\rm eff}}\left(\frac{G\mathcal M}{c^2}\right)^{5/3}
  \left(\frac{\pi f_{GW}(t)}{c}\right)^{2/3},
  \label{eq:htdamp}
\end{equation}
\emph{increases} with frequency, as a real inspiral must. An earlier version
of Case A's \code{run.py} used the frequency-\emph{domain}
stationary-phase-approximation scaling $|h(f)|\propto f^{-7/6}$ (which
decreases) as if it were this time-domain envelope --- caught by eye (the
plotted chirp amplitude was shrinking toward merger) and fixed; both scalings
are now separate, clearly labelled functions
(\code{restricted\_pn\_amplitude\_td} vs.\ \code{\_fd}) with a regression
check in \code{test\_chirp.py} that Eq.~\eqref{eq:htdamp} increases with
$f$.
No absolute strain calibration (inclination/polarization prefactors, folded
into the leading $4$) is claimed anywhere in this project --- only relative
amplitudes and the lensing amplification factor matter for the physics
studied here.

% ======================================================================
\chapter{Putting it together: two regimes}
\label{ch:regimes}
% ======================================================================

\section{Static lens (Case A)}
The lens is frozen for the whole chirp when the merger duration is much
shorter than the outer period, $\tau(f_{\rm start})\ll P_{\rm out}$. For the
parameters in \code{src/gwlens/system.py} ($f_{\rm start}=10$\,Hz,
$P_{\rm out}=4$\,d): $\tau=15.2$\,s, a fraction $4.4\times10^{-5}$ of
$P_{\rm out}$ --- static to extremely high accuracy
(\code{tests/test\_system.py::check\_static\_lens\_regime}). Lensing is then
a single frequency-domain multiplication, $h_{\rm lensed}(f)=F(f,y_A)\,
h_{\rm unlensed}(f)$, applied via FFT/IFFT in \code{cases/case\_A\_chirp/
run.py}.

\section{Quasi-monochromatic, moving lens (Case B)}
\label{sec:caseB}
Conversely, the lens genuinely moves when the observation spans several
outer periods while the GW frequency itself barely evolves: $\tau(f_B)\gg
T_{\rm obs}\gg P_{\rm out}$. For $f_B=0.05$\,Hz, $T_{\rm obs}=6P_{\rm out}$
(24\,d): $\tau=240.7$\,d ($60.2$ outer periods remain), and the fractional
frequency drift over $T_{\rm obs}$ is $4.0\%$
(\code{tests/test\_system.py::check\_quasi\_monochromatic\_regime}).

\subsection*{Why instantaneous (adiabatic) $F(w_B,y(t))$ is the right
prescription}
Eq.~\eqref{eq:Fdim}'s frequency-domain multiplication assumes a
\emph{fixed} lens configuration; it does not directly apply when $y$ itself
is a function of time. The separation of scales here is enormous: $y(t)$
varies on the outer-orbital timescale ($P_{\rm out}\sim$ days), while the GW
itself oscillates on a timescale $1/f_B\sim20$\,s --- nine orders of
magnitude faster. To the wave, the lens configuration is, to superb
approximation, constant over any one GW cycle (and over very many
consecutive cycles). This is the standard adiabatic/quasi-static
approximation: at each instant, apply the amplification factor the lens
\emph{would} produce for a static source at that instant's $y(t)$,
\begin{equation}
  z_{\rm lensed}(t) \;=\; F\!\big(w_B,\,y(t)\big)\;z_{\rm unlensed}(t),
  \qquad z(t)\equiv h(t)\,e^{i\phi(t)}\ \text{(analytic signal)},
  \label{eq:adiabatic}
\end{equation}
used in \code{cases/case\_B\_monochromatic/run.py}. This is the one place
this project's two cases use genuinely different numerical machinery (FFT
convolution for Case A, direct time-domain multiplication for Case B), and
the reason is this timescale argument, not a difference in the underlying
physics.

\section{An idealized detector}
Both cases stop at an idealized, illustrative detector view: the
analytic-signal (Hilbert-transform) envelope of $h_{\rm lensed}(t)$ vs.\
$h_{\rm unlensed}(t)$, with no noise curve, no antenna pattern, and no
matched filtering. This is a deliberate scope choice
(\code{wiki/conventions.md}): the point of this project is the wave-optics
lensing physics and its visualization, not a detection-pipeline study,
and the figures are explicit about the omission on their face.

% ======================================================================
\chapter*{References}
\addcontentsline{toc}{chapter}{References}
\bibliographystyle{plainnat}
\bibliography{refs}

\end{document}
